\documentclass[UserManual.tex]{subfiles}
\begin{document}
\begin{sloppypar}
\section{Binding {\it Smooth Emulator} to Python Code}\label{sec:python}

\subsection{Overview}
The User might wish to call the emulator functionality from within a Python script or from within a Python notebook. The ground work, based on pybind11, exists for calling some of the emulator functionality within a Python script. The User may create an object from within Python from which they can create and tune an emulator. Simple calls exist which permit the User to get emulator predictions and uncertainties as a function of either the model parameters, $\vec{X}$, or the scaled parameter, $\vec{\theta}$. Functions that translate back and forth between $\vec{X}$ and $\vec{\theta}$ are also provided.

The training-point-optimization and MCMC functionality is not available in the current software. However, this may be sufficient if a User wishes to build and tune an emulator using {\it Smooth Emulator} software, then use their own MCMC python-based software to explore the model-parameter space.

\subsection{Setting Up The Software}
The binding of C++ libraries to Python is accomplished through pybind11. Due to compatibility problems with pybind11, {\it Smooth Emulator} software is compiled with version 15 of {\tt g++}. 

\subsection{Editing and Running a Python Script}
An example Python script, {\tt smoothy\_emulate.py}, is in the \textit{SMOOTH\_HOME}{\tt /bin} directory. 

To run the script,
\vspace*{-8pt}{\tt
\begin{Verbatim}[commandchars=\\\{\}]
\textit{MY_ANALYSIS}% \textcolor{darkred}{\textit{SMOOTH_HOME}/bin/smoothy_emulate.py}
Adding this to path: XXXXXX
X= [20. 20. 20. 20. 20. 20.]
Y[ 0 ]= 62.06376520471999  SigmaY= 0.46481585069318604
Y[ 1 ]= -88.61735278099212  SigmaY= 2.198210063746222
Y[ 2 ]= 35.35391376800502  SigmaY= 0.5431950258044482
Y[ 3 ]= 47.38398026724699  SigmaY= 0.5632478557065526
Y[ 4 ]= 41.16532652137935  SigmaY= 0.29387511492428314
Y[ 5 ]= 64.91706104229243  SigmaY= 0.23423530131223003
Y[ 6 ]= -192.62294658870263  SigmaY= 0.04820162692778329
Y[ 7 ]= 88.43393115018455  SigmaY= 0.3480224686688987
Y[ 8 ]= 3.840354622268825  SigmaY= 0.5590394534630352
Y[ 9 ]= 64.70923479158401  SigmaY= 0.7742206531639135
\end{Verbatim}
}\vspace*{-8pt}
The script took the point in model-parameter space, $(20,20,20,20,20,20)$, and calculated all three observables along with their uncertainties.

The user might find it easier to move the Python script to the current directory so that they can avoid typing in the long path. The Python script, {\tt smoothy\_emulate.py}, is
\vspace*{-8pt}{\tt
\begin{Verbatim}[commandchars=\\\{\}]
import numpy as np
import sys
import os
print('Adding this to path: '+os.environ['SMOOTH_HOME'])
sys.path.insert(0,os.environ['SMOOTH_HOME']+"/software/pybind_stuff")
import emulator_smooth
smoothmaster=emulator_smooth.emulator_smooth()
smoothmaster.TuneAllY()
NPars=smoothmaster.GetNPars()
NObs=smoothmaster.GetNObs()

X=np.zeros(NPars,dtype='float')
Y=np.zeros(NObs,dtype='float')
SigmaY=np.zeros(NObs,dtype='float')

X[0]=20
X[1]=20
X[2]=20
X[3]=20
X[4]=20
X[5]=20
print('X=',X)

for iY in range(0,NObs):
  Y[iY],SigmaY[iY]=smoothmaster.GetYSigmaFromXPython(iY,X)
  print('Y[',iY,']=',Y[iY],' SigmaY=',SigmaY[iY])
\end{Verbatim}
}\vspace*{-8pt}
There are other calls available with the current software. One could add some version of the lines below to their Python script.
\vspace*{-8pt}{\tt
\begin{Verbatim}[commandchars=\\\{\}]
theta=smoothmaster.GetThetaFromX(X)
X=smoothmaster.GetXFromTheta(theta)
Y[iY],SigmaY[iY]=smoothmaster.GetYSigmaFromThetaPython(iY,theta)
Y[iY]=smoothmaster.GetYOnlyFromXPython(iY,X)
Y[iY]=smoothmaster.GetYOnlyFromThetaPython(iY,theta)
\end{Verbatim}
}\vspace*{-8pt}
The vectors {\tt X[..]} reference the model parameters, whereas {\tt theta[..]} describe the scaled model parameters. The emulated observables and the corresponding emulator uncertainties are {\tt Y[..]} and {\tt SigmaY[..]}. Hopefully, the User will find the functionality of this code to be self-explanatory. These calls should be sufficient to incorporate training and calling the emulator into any Python script.

\subsection{Making More Functionality Available via Python}
The easiest way to add more functions is to edit and extend the file \textit{SMOOTH\_HOME}{\tt /software/pybind\_stuff/pybind\_main.cc}. One can see how all the functions above were defined. One can extend this to more of the functions defined in the header file \textit{SMOOTH\_HOME}{\tt /software/include/msu\_smooth/master.h}. However, the syntax is rather opaque and not for the faint of heart. One should remember that the binding cannot easily change the functionality of functions with arguments passed by reference. Thus, if the User wishes to find some variable, e.g. {\tt Z}, from some C++ function that is written in the form {\tt void ThisIsMyMethodToGetZ(double \&Z,double x,double y)}, they will instead have to write a new C++ function in the form {\tt double ThisIsMyMethodToGetZ(double x,double z)}, where the new function returns {\tt Z} rather than having it updated by the function. For example to return the emulated value and uncertainty, a function was created of the form {\tt vector<double> GetYSigmaFromXPython(int iY,vector<double> X)}. The returned vector was 2-dimensional with {\tt Y[iY]} being the first element and {\tt sigmaY[iY]} being the second element.

\end{sloppypar}
\end{document}